\chapter{วิธีการใช้งานและการตั้งค่าเบื้องต้น}

\section{ความต้องการขั้นต่ำของชุดคำสั่ง}

\begin{enumerate}
  \item \lualatexlogo รุ่น (เวอร์ชัน) หลังเดือนตุลาคม ปี ค.ศ. 2025 (พ.ศ. 2568)
  \item ชุดคำสั่ง \texttt{fontspec} และ \texttt{polyglossia} ของ \latexlogo
  \item ชุดแบบอักษรไทย (Thai font)
\end{enumerate}

\paragraph{หมายเหตุ :}

ณ ขณะนี้ ยังไม่สามารถใช้ชุดคำสั่งนี้กับชุดคำสั่ง \texttt{babel} ได้

\section{วิธีการดาวน์โหลดชุดคำสั่งจากแหล่งต้นทาง}

\paragraph{ในระบบปฎิบัติการ Linux :}

ให้รันคำสั่งต่อไปนี้ใน bash

\begin{bashsource}{}
mkdir -p $(kpsewhich -var-value TEXMFHOME)/tex/latex
cd $(kpsewhich -var-value TEXMFHOME)/tex/latex
git clone https://github.com/AtomThum/thai-segmentation
rm -rf .git
\end{bashsource}

สามารถตรวจสอบว่าชุดคำสั่งพร้อมใช้งานได้ด้วย
\begin{bashsource}{}
kpsewhich thai-segmentation.sty
\end{bashsource}

\paragraph{ในระบบปฎิบัติการ Windows :}

\begin{enumerate}
  \item ถ้าไม่ได้ตั้งค่า \texttt{TEXMFHOME} ไว้ ให้สร้างโฟลเดอร์ \\
        \texttt{\%USERPROFILE\%\textbackslash texmf\textbackslash
          tex\textbackslash latex\textbackslash thai-segmentation}
  \item ไปที่ยูอาร์แอล \url{https://github.com/AtomThum/thai-segmentation}
  \item กดที่ปุ่ม \texttt{Code} สีเขียว แล้วกดดาวน์โหลด ZIP ไฟล์
  \item แตกไฟล์ \texttt{thai-segmentation-main.zip}
  \item คัดลอกไฟล์ในโฟลเดอร์ที่แตกออกมาไปที่ \\
        \texttt{\%USERPROFILE\%\textbackslash tex\textbackslash
          latex\textbackslash thai-segmentation}
\end{enumerate}

\paragraph{ในระบบปฎิบัติการ Macs :}

ยังไม่มีข้อมูลเนื่องจากผู้พัฒนาชุดคำสั่งไม่มีเงินเพื่อซื้อ Macbook

\section{การเรียกใช้งานชุดคำสั่งและการตั้งค่าเอกสาร}
\label{sec:configuration}

ตั้งเอกสารให้ใช้ \lualatexlogo และตั้งให้ \latexlogo ใช้ภาษาไทยโดยใช้
\texttt{polyglossia} เลือกชุดแบบอักษรที่ต้องการโดยใช \texttt{fontspec}
ตัวอย่างเช่น

\begin{texsource}{}
% !TEX program = lualatex
\documentclass{article}

\usepackage{fontspec}
\usepackage{polyglossia}

\newfontfamily{\thaifont}{Sarabun}[Script=Thai]
\newfontfamily{\thaifontsf}{Athiti}[Script=Thai,Scale=MatchUppercase]
\newfontfamily{\englishfont}{Sarabun}[Ligatures=TeX]

\setdefaultlanguage{thai}
\setmainlanguage{thai}
\setotherlanguage{english}
\end{texsource}

ณ ขณะนี้ ยังไม่สามารถใช้ชุดคำสั่งนี้กับชุดคำสั่ง \texttt{babel} ได้

สามารถเรียกใช้ชุดคำสั่งจาก \texttt{thai-segmentation} ได้โดย

\begin{texsource}{}
  \usepackage[<options>]{thai-segmentation}
\end{texsource}
สามารถปรับค่าต่าง ๆ ใน \texttt{<options>} ได้ดังนี้
\begin{description}[font=\ttfamily]
  \item[thai-segmentation] (\texttt{true}/\texttt{false}) เปิด-ปิดการตัดคำไทย
  \item[interword-expansion] (\texttt{true}/\texttt{false})
        เปิด-ปิดการขยายช่องว่างระหว่างคำในประโยค
        (เพิ่มเติมใน\zcref{sec:interword-expansion})
  \item[interword-stretch] (ทศนิยม) จะมีผลก็ต่อเมื่อ
        \texttt{interword-expansion=true} เท่านั้น
        กำหนดขนาดที่ช่องว่างระหว่างคำสามารถขยายได้ มีหน่วยเป็นพอยต์ (pt)
        มีค่าโดยปริยายเป็น \texttt{1} ไม่ควรตั้งค่าเกินกว่านี้
        เนื่องจากอาจทำให้คำในประโยคห่างกันมากเกินไป
  \item[interword-shrink] (ทศนิยม) จะมีผลก็ต่อเมื่อ
        \texttt{interword-expansion=true} เท่านั้น
        กำหนดขนาดสูงสุดที่ช่องหว่างระหว่างคำสามารถหดได้ มีหน่วยเป็นพอยต์ (pt)
        มีค่าโดยปริยายเป็น \texttt{0.5}
        ไม่ควรตั้งค่าเกินกว่านี้เนื่องจากอาจทำให้คำชิดกันเกินไปหรือซ้อนทับกัน
  \item[intercluster-stretch] (ทศนิยม)
        กำหนดระยะห่างเพิ่มเติมระหว่างกลุ่มตัวอักษรไทย (Thai character clusters)
        จะมีผลก็ต่อเมื่อตั้งเป็นค่าทศนิยมบวก มีค่าปริยายเป็น \texttt{0}
        ควรตั้งค่าให้น้อยกว่า \texttt{interword-stretch}
  \item[dictionary] ระบุที่อยู่ (path) ของไฟล์คำที่จะใส่เพิ่มเข้าพจนานุกรมที่
        \texttt{newmm} ใช้ตัดคำ ตัวอย่างเช่น
        ถ้าโครงสร้างโฟลเดอร์เป็นแบบ\zcref{fig:folder-example} โดยไฟล์
        \texttt{file.txt} มีเนื้อหาดังนี้
        \begin{texsource}{file.txt}
ไอน์สไตน์
พลังค์
แมกซ์เวลล์
\end{texsource}
        สามารถเพิ่มคำจากไฟล์ \texttt{file.txt} เข้าพจนานุกรมได้โดยตั้งค่า \\
        \texttt{dictionaries=file.txt}
        (เพิ่มเติมใน\zcref{sec:interword-expansion})
        \begin{figure}[h]
          \centering
          \begin{forest}
            dir tree [project/ [file.txt] [main.tex] [main.pdf] ]
          \end{forest}
          \caption{ตัวอย่างโครงสร้างโฟลเดอร์}
          \label{fig:folder-example}
        \end{figure}
  \item[tcc-min-edge] (จำนวนเต็ม) หากคำใดมีตัวอักษรน้อยกว่าค่านี้
        จะไม่ใส่จุดตัดคำด้วยยัติภังค์ มีค่าโดยปริยายเป็น \texttt{2}
        (เพิ่มเติมใน\zcref{sec:interword-expansion})
  \item[show-hyph] (\texttt{true}/\texttt{false})
        เป็นเครื่องมือสำหรับผู้พัฒนาเพื่อตรวจสอบจุดตัดคำและตำแหน่งการใส่ยัติภังค์
        ไม่ควรเปิดใช้ในเอกสารฉบับสมบูรณ์
\end{description}

ตัวอย่างการตั้งค่าและเรียกใช้ชุดคำสั่ง เช่น
\begin{texsource}{}
\usepackage[
  interword-expansion=true,
  interword-stretch=0.5,
  interword-shrink=0.2,
  intercluster-stretch=0,
  dictionary=dict.txt,
  tcc-min-edge=3,
]{thai-segmentation}
\end{texsource}

\section{รายชื่อคำสั่ง}

\begin{description}[font=\ttfamily]
  \item[\textbackslash ThaiAddWords\{<คำที่หนึ่ง,คำที่สอง,...>\}] เพิ่มคำเข้าพจนานุกรมของ
        \texttt{newmm} รับค่าเหมือน \texttt{dictionaries} ในการตั้งค่าชุดคำสั่ง
        ตัวอย่างวิธีการใช้เช่น \\
        \texttt{\textbackslash ThaiAddwords\{ลากรานจ์,แฮมิลตัน,แฮมิลโทเนียน\}}
  \item[\textbackslash ThaiRemoveWords\{<คำที่หนึ่ง,คำที่สอง,...>\}] ลบคำจากพจนานุกรมของ
        \texttt{newmm} รับค่าเหมือน \texttt{dictionaries} ในการตั้งค่าชุดคำสั่ง
        ตัวอย่างการใช้เช่น \\
        \texttt{\textbackslash ThaiRemoveWords\{ยั่งยืน,อย่างง่าย\}}
  \item[\textbackslash ThaiAddHyphenation\{<คำ-ที่-หนึ่ง,คำ-ที่-สอง,...>\}] กำหนดตำแหน่งการใส่ยัติภังค์
        กรณีที่ระบบตัดคำใส่จุดตัดคำที่ไม่ตรงกับการออกเสียงหรือคำอ่านธรรมชาติ
        เช่น หากระบบตัดคำว่า "ลากรานจ์" เป็น "ลาก-รานจ์"
        ซึ่งไม่ตรงกับคำอ่านธรรมชาติ สามารถแก้ไขให้ถูกต้องได้ด้วยคำสั่ง \\
        \texttt{\textbackslash ThaiAddHyphenation\{ลา-กรานจ์\}}
  \item[\textbackslash ThaiRemoveHyphenation\{คำ-ที่-หนึ่ง,คำ-ที่-สอง\}] ใช้เหมือน
        \texttt{\textbackslash ThaiAddHyphenation}
        แต่ใช้เพื่อลบคำสั่งบังคับยัติภังค์ที่กำหนดไว้ก่อนหน้า
  \item[\textbackslash ThaiAddWordsFromFile\{<ไฟล์พจนานุกรม>\}] นำเข้าไฟล์พจนานุกรมเพิ่มเติม
        การทำงานเหมือนการกำหนดผ่านการตั้งค่า \texttt{dictionary}
        (\zcref{sec:configuration})
  \item[\textbackslash ThaiAddSyllables\{<กลุ่มอักษรที่หนึ่ง,กลุ่มอักษรที่สอง,...>\}] เพิ่มกลุ่มอักษรเพื่อใช้ในการใส่ยัติภังค์
  \item[\textbackslash ThaiRemoveSyllables\{<กลุ่มอักษรที่หนึ่ง,กลุ่มอักษรที่สอง>\}] ลบกลุ่มอักษรออกจากระบบการใส่ยัติภังค์
\end{description}
